\section{Giới thiệu}

Khi số chiều đặc trưng rất lớn, nhiều thuật toán học máy tốn kém bộ nhớ và thời gian, đồng thời khó quan sát cấu trúc dữ liệu. \emph{Giảm chiều} (dimensionality reduction) nhằm tìm một biểu diễn có số chiều nhỏ hơn nhưng vẫn giữ được phần thông tin hoặc quan hệ hình học quan trọng của tập mẫu gốc.

\paragraph{Quy ước ký hiệu.}
\begin{itemize}
  \item \textbf{Tập hợp:} $\mathcal{X}$ (không gian mẫu), $\mathcal{H}$ (không gian Hilbert), $P_k$ (tập ma trận chiếu).
  \item \textbf{Vector:} chữ thường in đậm, ví dụ $\mathbf{x}_i \in \R^N$.
  \item \textbf{Ma trận:} chữ hoa in đậm, ví dụ $\mathbf{X} \in \R^{N \times m}$.
  \item \textbf{Vô hướng (scalar):} chữ nghiêng thường, ví dụ $N$, $m$, $k$, $\varepsilon$, $\lambda$.
\end{itemize}

\subsection{Bài toán và ký hiệu}

Cho tập mẫu $S = (x_1, \ldots, x_m)$ với $x_i \in \mathcal{X}$, ánh xạ đặc trưng $\Phi : \mathcal{X} \to \R^N$ và ma trận dữ liệu
\[
\mathbf{X} = \big[\Phi(x_1), \ldots, \Phi(x_m)\big] \in \R^{N \times m},
\]
cột thứ $i$ là vector $\mathbf{x}_i = \Phi(x_i)$. Ta muốn tìm ma trận biểu diễn $\mathbf{Y} \in \R^{k \times m}$ với $k \ll N$ sao cho $\mathbf{Y}$ phản ánh $\mathbf{X}$ theo tiêu chí đã chọn (lỗi tái tạo, phương sai, khoảng cách, láng giềng,\ldots).

\subsection{Động lực thực hiện}

\begin{itemize}
  \item \textbf{Tính toán:} nén dữ liệu trước khi huấn luyện hoặc truy vấn, giảm chi phí lưu trữ và phép toán.
  \item \textbf{Trực quan:} đưa mẫu về hai hoặc ba chiều để khám phá cụm, ngoại lệ, xu hướng.
  \item \textbf{Trích đặc trưng:} tạo tập chiều thấp hy vọng hữu ích hơn cho phân loại hoặc hồi quy phía sau.
\end{itemize}

\subsection{Ví dụ Swiss roll}

Swiss roll là dữ liệu mô phỏng ba chiều nằm gần một mặt hai chiều ``cuộn'' (Hình~\ref{fig:swiss-roll}). Nếu ``mở'' mặt này ra phẳng, quan hệ láng giềng trở nên rõ. Trong thực tế, manifold lý tưởng hiếm khi khớp hoàn toàn; ví dụ này chủ yếu giúp hình dung vì sao cần phương pháp \emph{phi tuyến}, chứ không phải bằng chứng cho mọi bộ dữ liệu thật.

\begin{figure}[H]
\centering
\begin{tikzpicture}[scale=0.9, every node/.style={font=\small}]
  \begin{scope}
    \node[font=\bfseries] at (0,2.6) {(a) Dữ liệu ``cuộn'' (minh họa)};
    \draw[blue!60, thick, domain=0:720, samples=120, smooth]
      plot ({(0.8+0.003*\x)*cos(\x)}, {(0.8+0.003*\x)*sin(\x)});
    \foreach \t in {60,180,300,480,600} {
      \fill[red!70] ({(0.8+0.003*\t)*cos(\t)}, {(0.8+0.003*\t)*sin(\t)}) circle (1.8pt);
    }
  \end{scope}
  \begin{scope}[xshift=5.5cm]
    \node[font=\bfseries] at (0,2.6) {(b) Mở ra mặt phẳng};
    \draw[fill=blue!8, draw=blue!30] (-2.1,-0.6) rectangle (2.1,2.2);
    \foreach \x/\y in {-1.4/1.6, -0.4/0.8, 0.6/1.4, 1.3/0.5, -0.9/0.3, 1.0/1.9} {
      \fill[blue!70] (\x,\y) circle (2pt);
    }
  \end{scope}
\end{tikzpicture}
\caption{Ý tưởng Swiss roll: mẫu nằm trên cấu trúc hai chiều được nhúng trong không gian cao chiều; sau giảm chiều, quan hệ láng giềng dễ quan sát hơn.}
\label{fig:swiss-roll}
\end{figure}

\subsection{Đóng góp của báo cáo}

\begin{enumerate}
  \item Chương~2: cơ sở lý thuyết---đại số ma trận (hiệp phương sai, SVD, chiếu, kernel, Laplacian), hình học manifold, phân kỳ KL / cross-entropy, và khung ước lượng tương phản NCE / NEG.
  \item Chương~3: các phương pháp giảm chiều theo \cite{mohri2018foundations}---PCA, KPCA, Isomap, Laplacian eigenmaps, LLE, Johnson--Lindenstrauss; kèm \textbf{[MỞ RỘNG]} t-SNE và UMAP \cite{vandermaaten2008tsne,mcinnes2018umap} (không có trong sách).
  \item Chương~4: khảo sát bài báo gần đây---thống nhất t-SNE và UMAP qua học tương phản \cite{damrich2023contrastive}.
  \item Chương~5: thực nghiệm (đang bổ sung).
  \item Chương~6: tổng kết.
  \item Phụ lục: các chứng minh kỹ thuật hỗ trợ (bổ đề tập trung $\chi^2$, khai triển gradient UMAP và NCE).
\end{enumerate}

Nguồn tham khảo chính cho phần lý thuyết theo sách là \emph{Foundations of Machine Learning} \cite{mohri2018foundations}; t-SNE và UMAP trích từ bài báo tương ứng.
